% Options for packages loaded elsewhere
\PassOptionsToPackage{unicode}{hyperref}
\PassOptionsToPackage{hyphens}{url}
\documentclass[
  mongolian,
]{article}
\usepackage{xcolor}
\usepackage{amsmath,amssymb}
\setcounter{secnumdepth}{-\maxdimen} % remove section numbering
\usepackage{iftex}
\ifPDFTeX
  \usepackage[T1]{fontenc}
  \usepackage[utf8]{inputenc}
  \usepackage{textcomp} % provide euro and other symbols
\else % if luatex or xetex
  \usepackage{unicode-math} % this also loads fontspec
  \defaultfontfeatures{Scale=MatchLowercase}
  \defaultfontfeatures[\rmfamily]{Ligatures=TeX,Scale=1}
\fi
\usepackage{lmodern}
\ifPDFTeX\else
  % xetex/luatex font selection
\fi
% Use upquote if available, for straight quotes in verbatim environments
\IfFileExists{upquote.sty}{\usepackage{upquote}}{}
\IfFileExists{microtype.sty}{% use microtype if available
  \usepackage[]{microtype}
  \UseMicrotypeSet[protrusion]{basicmath} % disable protrusion for tt fonts
}{}
\makeatletter
\@ifundefined{KOMAClassName}{% if non-KOMA class
  \IfFileExists{parskip.sty}{%
    \usepackage{parskip}
  }{% else
    \setlength{\parindent}{0pt}
    \setlength{\parskip}{6pt plus 2pt minus 1pt}}
}{% if KOMA class
  \KOMAoptions{parskip=half}}
\makeatother
\usepackage{color}
\usepackage{fancyvrb}
\newcommand{\VerbBar}{|}
\newcommand{\VERB}{\Verb[commandchars=\\\{\}]}
\DefineVerbatimEnvironment{Highlighting}{Verbatim}{commandchars=\\\{\}}
% Add ',fontsize=\small' for more characters per line
\newenvironment{Shaded}{}{}
\newcommand{\AlertTok}[1]{\textcolor[rgb]{1.00,0.00,0.00}{\textbf{#1}}}
\newcommand{\AnnotationTok}[1]{\textcolor[rgb]{0.38,0.63,0.69}{\textbf{\textit{#1}}}}
\newcommand{\AttributeTok}[1]{\textcolor[rgb]{0.49,0.56,0.16}{#1}}
\newcommand{\BaseNTok}[1]{\textcolor[rgb]{0.25,0.63,0.44}{#1}}
\newcommand{\BuiltInTok}[1]{\textcolor[rgb]{0.00,0.50,0.00}{#1}}
\newcommand{\CharTok}[1]{\textcolor[rgb]{0.25,0.44,0.63}{#1}}
\newcommand{\CommentTok}[1]{\textcolor[rgb]{0.38,0.63,0.69}{\textit{#1}}}
\newcommand{\CommentVarTok}[1]{\textcolor[rgb]{0.38,0.63,0.69}{\textbf{\textit{#1}}}}
\newcommand{\ConstantTok}[1]{\textcolor[rgb]{0.53,0.00,0.00}{#1}}
\newcommand{\ControlFlowTok}[1]{\textcolor[rgb]{0.00,0.44,0.13}{\textbf{#1}}}
\newcommand{\DataTypeTok}[1]{\textcolor[rgb]{0.56,0.13,0.00}{#1}}
\newcommand{\DecValTok}[1]{\textcolor[rgb]{0.25,0.63,0.44}{#1}}
\newcommand{\DocumentationTok}[1]{\textcolor[rgb]{0.73,0.13,0.13}{\textit{#1}}}
\newcommand{\ErrorTok}[1]{\textcolor[rgb]{1.00,0.00,0.00}{\textbf{#1}}}
\newcommand{\ExtensionTok}[1]{#1}
\newcommand{\FloatTok}[1]{\textcolor[rgb]{0.25,0.63,0.44}{#1}}
\newcommand{\FunctionTok}[1]{\textcolor[rgb]{0.02,0.16,0.49}{#1}}
\newcommand{\ImportTok}[1]{\textcolor[rgb]{0.00,0.50,0.00}{\textbf{#1}}}
\newcommand{\InformationTok}[1]{\textcolor[rgb]{0.38,0.63,0.69}{\textbf{\textit{#1}}}}
\newcommand{\KeywordTok}[1]{\textcolor[rgb]{0.00,0.44,0.13}{\textbf{#1}}}
\newcommand{\NormalTok}[1]{#1}
\newcommand{\OperatorTok}[1]{\textcolor[rgb]{0.40,0.40,0.40}{#1}}
\newcommand{\OtherTok}[1]{\textcolor[rgb]{0.00,0.44,0.13}{#1}}
\newcommand{\PreprocessorTok}[1]{\textcolor[rgb]{0.74,0.48,0.00}{#1}}
\newcommand{\RegionMarkerTok}[1]{#1}
\newcommand{\SpecialCharTok}[1]{\textcolor[rgb]{0.25,0.44,0.63}{#1}}
\newcommand{\SpecialStringTok}[1]{\textcolor[rgb]{0.73,0.40,0.53}{#1}}
\newcommand{\StringTok}[1]{\textcolor[rgb]{0.25,0.44,0.63}{#1}}
\newcommand{\VariableTok}[1]{\textcolor[rgb]{0.10,0.09,0.49}{#1}}
\newcommand{\VerbatimStringTok}[1]{\textcolor[rgb]{0.25,0.44,0.63}{#1}}
\newcommand{\WarningTok}[1]{\textcolor[rgb]{0.38,0.63,0.69}{\textbf{\textit{#1}}}}
\usepackage{longtable,booktabs,array}
\usepackage{calc} % for calculating minipage widths
% Correct order of tables after \paragraph or \subparagraph
\usepackage{etoolbox}
\makeatletter
\patchcmd\longtable{\par}{\if@noskipsec\mbox{}\fi\par}{}{}
\makeatother
% Allow footnotes in longtable head/foot
\IfFileExists{footnotehyper.sty}{\usepackage{footnotehyper}}{\usepackage{footnote}}
\makesavenoteenv{longtable}
\ifLuaTeX
\usepackage[bidi=basic,shorthands=off,,provide=*]{babel}
\else
\usepackage[bidi=default,shorthands=off,,provide=*]{babel}
\fi
\ifLuaTeX
  \usepackage{selnolig} % disable illegal ligatures
\fi
\setlength{\emergencystretch}{3em} % prevent overfull lines
\providecommand{\tightlist}{%
  \setlength{\itemsep}{0pt}\setlength{\parskip}{0pt}}
\usepackage{fontspec}
\setmainfont{Liberation Serif}
\setsansfont{Liberation Sans}
\setmonofont{Liberation Mono}[Scale=0.86]
\usepackage[a4paper,left=3cm,right=2cm,top=2cm,bottom=2cm,includefoot,headheight=15pt]{geometry}
\usepackage{ragged2e,indentfirst,xurl}
\usepackage{longtable,booktabs,array,calc,graphicx}
\usepackage{xcolor,fancyhdr,titlesec,enumitem,float}
\definecolor{ink}{HTML}{19354A}
\definecolor{accent}{HTML}{167D8D}
\titleformat{\section}{\Large\sffamily\bfseries\color{ink}}{}{0pt}{}
\titleformat{\subsection}{\large\sffamily\bfseries\color{ink}}{}{0pt}{}
\titleformat{\subsubsection}{\fontsize{11.5}{13.5}\selectfont\sffamily\bfseries\color{ink}}{}{0pt}{}
\titlespacing*{\section}{1.25cm}{0pt}{8pt}
\titlespacing*{\subsection}{1.25cm}{8pt}{4pt}
\titlespacing*{\subsubsection}{1.25cm}{6pt}{3pt}
\setlength{\parindent}{1.25cm}
\setlength{\parskip}{3pt}
\setlength{\emergencystretch}{3em}
\setlength{\tabcolsep}{4pt}
\renewcommand{\arraystretch}{1.08}
\setlist{nosep,leftmargin=1.5em,topsep=3pt}
\pagestyle{fancy}\fancyhf{}
\fancyhead[L]{\small\sffamily ICSI438 / SEMINAR 05}
\fancyhead[R]{\small\sffamily Nogoolin / OpenAPI}
\fancyfoot[R]{\small\thepage}
\renewcommand{\headrulewidth}{0.3pt}
\AtBeginDocument{\fontsize{10.5}{12.5}\selectfont\justifying\setlength{\parindent}{1.25cm}}


\AtBeginDocument{\renewcommand{\figurename}{Зураг}}
\usepackage{bookmark}
\IfFileExists{xurl.sty}{\usepackage{xurl}}{} % add URL line breaks if available
\urlstyle{same}
\hypersetup{
  pdflang={mn},
  hidelinks,
  pdfcreator={LaTeX via pandoc}}

\author{}
\date{}

\begin{document}

\hypersetup{
  pdftitle={Програм хангамжийн төслийн баримт бичиг — Seminar 5 OpenAPI},
  pdfauthor={Ц. Тэнгис}
}
\pagenumbering{gobble}
\thispagestyle{empty}
\noindent
\begin{minipage}[t][\dimexpr\textheight-\topskip\relax][t]{\linewidth}
  \begin{center}
    \textbf{МОНГОЛ УЛСЫН ИХ СУРГУУЛЬ}\\
    \textbf{МЭДЭЭЛЛИЙН ТЕХНОЛОГИ, ЭЛЕКТРОНИКИЙН СУРГУУЛЬ}\\
    \textbf{МЭДЭЭЛЭЛ, КОМПЬЮТЕРЫН УХААНЫ ТЭНХИМ}\\[1.7cm]
    \includegraphics[width=6cm]{../assets/num-logo.jpg}\\[1.2cm]
    {\fontsize{16}{20}\selectfont\bfseries
      ПРОГРАМ ХАНГАМЖИЙН ТӨСЛИЙН БАРИМТ БИЧИГ\par}
    \vspace{0.45cm}
    {\large\bfseries SEMINAR 5 — OPENAPI 3.0 БА CODE SAMPLE}
  \end{center}

  \vspace{1.2cm}
  \begin{tabular}{>{\raggedright\arraybackslash}m{5cm}@{\hspace{1cm}}>{\raggedright\arraybackslash}m{\dimexpr\linewidth-6cm-2\tabcolsep\relax}}
    Хичээл: & ICSI438 — Software Project Documentation \\[0.55cm]
    Үндсэн төсөл: & Nogoolin \\[0.55cm]
    Хичээлийн багш: & Ph.D. Batnyam Battulga \\[0.55cm]
    Гүйцэтгэсэн: & Програм хангамжийн IV ангийн оюутан\\
                  & Ц. Тэнгис (22B1NUM6249) \\[0.55cm]
    Баримтын хувилбар: & v0.1 draft
  \end{tabular}

  \vfill
  \begin{center}
    Улаанбаатар\\[0.25cm]
    2026 он
  \end{center}
\end{minipage}
\clearpage
\pagenumbering{arabic}
\setcounter{page}{1}


\section{Seminar 5 --- OpenAPI 3.0 ба code
sample}\label{seminar-5-openapi-3.0-ux431ux430-code-sample}

Төслийн GitHub дахь Sem5 хавтас:
\url{https://github.com/Tengis01/Nogoolin/tree/main/docs/ICSI438/sem5}.
Түүний \url{renderers/} хавтаст Redoc-ийн \texttt{redoc.html} болон
Swagger UI-ийн \texttt{swagger.html} байна.

\subsection{Зорилго ба
хүрээ}\label{ux437ux43eux440ux438ux43bux433ux43e-ux431ux430-ux445ux4afux440ux44dux44d}

Nogoolin-ийн одоо кодод бүртгэлтэй таван үйлдлийг
\href{openapi/openapi.yaml}{OpenAPI 3.0.3 YAML}-д тусгав. Сонголт нь
Sem3-ын inquiry SRS, Sem4-ийн inquiry runtime view-тэй шууд холбогдоно.
Энэ нь \texttt{docs/phase-0/06-api-spec.yaml}-ийн бүх API-г орлуулах
шинэ үндсэн spec биш; W5 лабораторийн нарийвчилсан slice юм. Серверийн
URL нь \texttt{localhost} тул Swagger-ийн ``Try it out'' зөвхөн local
API, local Supabase, тохирох auth token бэлэн үед ажиллана.

\begin{longtable}[]{@{}
  >{\raggedright\arraybackslash}p{(\linewidth - 8\tabcolsep) * \real{0.2000}}
  >{\raggedright\arraybackslash}p{(\linewidth - 8\tabcolsep) * \real{0.2000}}
  >{\raggedright\arraybackslash}p{(\linewidth - 8\tabcolsep) * \real{0.2000}}
  >{\raggedright\arraybackslash}p{(\linewidth - 8\tabcolsep) * \real{0.2000}}
  >{\raggedright\arraybackslash}p{(\linewidth - 8\tabcolsep) * \real{0.2000}}@{}}
\toprule\noalign{}
\begin{minipage}[b]{\linewidth}\raggedright
Үйлдэл
\end{minipage} & \begin{minipage}[b]{\linewidth}\raggedright
Sem3 шаардлага
\end{minipage} & \begin{minipage}[b]{\linewidth}\raggedright
Sem4 холбоо
\end{minipage} & \begin{minipage}[b]{\linewidth}\raggedright
Оролт
\end{minipage} & \begin{minipage}[b]{\linewidth}\raggedright
Амжилт
\end{minipage} \\
\midrule\noalign{}
\endhead
\bottomrule\noalign{}
\endlastfoot
\texttt{GET\ /products} & Каталогийн контекст & C4-02 API & query &
\texttt{200\ ProductList} \\
\texttt{POST\ /inquiries} & FR-01/02/03/05 & C4-D01 & JSON body; JWT
сонголтот & \texttt{201\ InquiryOne} \\
\texttt{GET\ /inquiries/mine} & FR-06, NFR-01 & C4-D01 & JWT &
\texttt{200\ InquiryArray} \\
\texttt{GET\ /admin/inquiries} & FR-07/08 & C4-D03 & query + admin JWT &
\texttt{200\ InquiryList} \\
\texttt{PATCH\ status} & FR-09 & C4-D03 & path + JSON body + admin JWT &
\texttt{200\ InquiryOne} \\
\end{longtable}

\texttt{PATCH\ status} нь \texttt{PATCH\ /admin/inquiries/\{id\}/status}
гэсэн бүтэн замтай. GET үйлдэлд \texttt{requestBody} бичээгүй. Оролтгүй
\texttt{GET\ /inquiries/mine} нь JWT-ээр хэрэглэгчээ танина; бусад
GET-ийн filter-ийг query parameter-аар өгнө. POST/PATCH-ийн жишээ болон
response schema YAML-д байна. \texttt{400/401/403/404/429/500} нь тухайн
үйлдэлд хэрэгтэй үед schema-тай. YAML дахь жишээ ID, огноо нь
\textbf{схемийг тайлбарлах өгөгдөл}; серверээс авсан JSON биш.
\texttt{product} холбоос response-д null эсвэл объект байх боломжтой;
кодын join-оор нягтална.

Нэг drift илэрсэн: Sem3 FR-03 inquiry утсыг 8 цифр гэж шаарддаг, харин
одоогийн \texttt{phoneSchema} нь \texttt{+}, зай, \texttt{-} тэмдэг
оролцсон 8--15 тэмдэгт зөвшөөрдөг. W5 YAML нь \textbf{одоогийн API
кодын} дүрмийг тусгасан; SRS ба кодын аль нэгийг багш/төслийн шийдвэрээр
дараа нь нэг мөр болгох хэрэгтэй.

\subsection{UE-5 --- Corg.ly Python
samples}\label{ue-5-corg.ly-python-samples}

Bhatti Ch.5 х.86--94-ийн контекст, товч байдал, нэршил, шууд ашиглалт,
бодит хариунд нийцэх гэсэн таван хэмжүүрээр гурван тусдаа Python script
бэлдэв. Эдгээр нь handout-ийн сургалтын \textbf{Corg.ly scenario};
Corg.ly-ийн ажилладаг сервер, батлагдсан request field contract, token
бидэнд байхгүй. Тиймээс Python синтаксийг шалгасан ч endpoint-уудыг live
ажиллуулсан, статус эсвэл returned JSON баталсан гэж үзэхгүй. Script нь
\texttt{CORGLY\_BASE\_URL}, \texttt{CORGLY\_TOKEN} environment variable
шаардаж, зөвхөн хэрэглэгч гараар ажиллуулбал HTTP хүсэлт явуулна.

\subsubsection{Зураг
илгээх}\label{ux437ux443ux440ux430ux433-ux438ux43bux433ux44dux44dux445}

\texttt{upload\_photo.py} нь бүртгэлтэй pet ID болон локал JPEG файлыг
multipart-аар илгээнэ. CLI-д \texttt{photo} ба \texttt{-\/-pet-id}
дамжуулна; \texttt{photo} field нэр нь handout-д тодорхойгүй тул
серверийн жинхэнэ contract-той тулгах шаардлагатай.

\begin{Shaded}
\begin{Highlighting}[]
\CommentTok{"""UE{-}5 Corg.ly exercise: upload one pet image. Contract is unverified."""}

\ImportTok{import}\NormalTok{ argparse}
\ImportTok{import}\NormalTok{ os}
\ImportTok{from}\NormalTok{ pathlib }\ImportTok{import}\NormalTok{ Path}

\ImportTok{import}\NormalTok{ requests}


\KeywordTok{def}\NormalTok{ main() }\OperatorTok{{-}\textgreater{}} \VariableTok{None}\NormalTok{:}
\NormalTok{    parser }\OperatorTok{=}\NormalTok{ argparse.ArgumentParser(description}\OperatorTok{=}\StringTok{"Upload a pet photo to a configured lab server"}\NormalTok{)}
\NormalTok{    parser.add\_argument(}\StringTok{"photo"}\NormalTok{, }\BuiltInTok{type}\OperatorTok{=}\NormalTok{Path, }\BuiltInTok{help}\OperatorTok{=}\StringTok{"Existing pet image path"}\NormalTok{)}
\NormalTok{    parser.add\_argument(}\StringTok{"{-}{-}pet{-}id"}\NormalTok{, required}\OperatorTok{=}\VariableTok{True}\NormalTok{, }\BuiltInTok{help}\OperatorTok{=}\StringTok{"Registered pet ID"}\NormalTok{)}
\NormalTok{    args }\OperatorTok{=}\NormalTok{ parser.parse\_args()}
\NormalTok{    base\_url }\OperatorTok{=}\NormalTok{ os.environ[}\StringTok{"CORGLY\_BASE\_URL"}\NormalTok{].rstrip(}\StringTok{"/"}\NormalTok{)}
\NormalTok{    token }\OperatorTok{=}\NormalTok{ os.environ[}\StringTok{"CORGLY\_TOKEN"}\NormalTok{]}
    \ControlFlowTok{with}\NormalTok{ args.photo.}\BuiltInTok{open}\NormalTok{(}\StringTok{"rb"}\NormalTok{) }\ImportTok{as}\NormalTok{ image:}
\NormalTok{        response }\OperatorTok{=}\NormalTok{ requests.post(}
            \SpecialStringTok{f"}\SpecialCharTok{\{}\NormalTok{base\_url}\SpecialCharTok{\}}\SpecialStringTok{/v1/pets/upload{-}photo"}\NormalTok{,}
\NormalTok{            headers}\OperatorTok{=}\NormalTok{\{}\StringTok{"Authorization"}\NormalTok{: }\SpecialStringTok{f"Bearer }\SpecialCharTok{\{}\NormalTok{token}\SpecialCharTok{\}}\SpecialStringTok{"}\NormalTok{\},}
\NormalTok{            data}\OperatorTok{=}\NormalTok{\{}\StringTok{"pet\_id"}\NormalTok{: args.pet\_id\},}
\NormalTok{            files}\OperatorTok{=}\NormalTok{\{}\StringTok{"photo"}\NormalTok{: (args.photo.name, image, }\StringTok{"image/jpeg"}\NormalTok{)\},}
\NormalTok{            timeout}\OperatorTok{=}\DecValTok{20}\NormalTok{,}
\NormalTok{        )}
    \BuiltInTok{print}\NormalTok{(response.status\_code, response.text)}
\NormalTok{    response.raise\_for\_status()}


\ControlFlowTok{if} \VariableTok{\_\_name\_\_} \OperatorTok{==} \StringTok{"\_\_main\_\_"}\NormalTok{:}
\NormalTok{    main()}
\end{Highlighting}
\end{Shaded}

\subsubsection{Bark
орчуулах}\label{bark-ux43eux440ux447ux443ux443ux43bux430ux445}

\texttt{translate\_bark.py} нь pet ID болон локал WAV бичлэгийг
multipart-аар илгээнэ. \texttt{audio} field нэр болон MIME нь
лабораторийн боломжит таамаг тул Corg.ly API-ийн баталгаа биш.

\begin{Shaded}
\begin{Highlighting}[]
\CommentTok{"""UE{-}5 Corg.ly exercise: request bark translation. Contract is unverified."""}

\ImportTok{import}\NormalTok{ argparse}
\ImportTok{import}\NormalTok{ os}
\ImportTok{from}\NormalTok{ pathlib }\ImportTok{import}\NormalTok{ Path}

\ImportTok{import}\NormalTok{ requests}


\KeywordTok{def}\NormalTok{ main() }\OperatorTok{{-}\textgreater{}} \VariableTok{None}\NormalTok{:}
\NormalTok{    parser }\OperatorTok{=}\NormalTok{ argparse.ArgumentParser(description}\OperatorTok{=}\StringTok{"Send a bark recording to a configured lab server"}\NormalTok{)}
\NormalTok{    parser.add\_argument(}\StringTok{"audio"}\NormalTok{, }\BuiltInTok{type}\OperatorTok{=}\NormalTok{Path, }\BuiltInTok{help}\OperatorTok{=}\StringTok{"Existing WAV recording path"}\NormalTok{)}
\NormalTok{    parser.add\_argument(}\StringTok{"{-}{-}pet{-}id"}\NormalTok{, required}\OperatorTok{=}\VariableTok{True}\NormalTok{, }\BuiltInTok{help}\OperatorTok{=}\StringTok{"Registered pet ID"}\NormalTok{)}
\NormalTok{    args }\OperatorTok{=}\NormalTok{ parser.parse\_args()}
\NormalTok{    base\_url }\OperatorTok{=}\NormalTok{ os.environ[}\StringTok{"CORGLY\_BASE\_URL"}\NormalTok{].rstrip(}\StringTok{"/"}\NormalTok{)}
\NormalTok{    token }\OperatorTok{=}\NormalTok{ os.environ[}\StringTok{"CORGLY\_TOKEN"}\NormalTok{]}
    \ControlFlowTok{with}\NormalTok{ args.audio.}\BuiltInTok{open}\NormalTok{(}\StringTok{"rb"}\NormalTok{) }\ImportTok{as}\NormalTok{ recording:}
\NormalTok{        response }\OperatorTok{=}\NormalTok{ requests.post(}
            \SpecialStringTok{f"}\SpecialCharTok{\{}\NormalTok{base\_url}\SpecialCharTok{\}}\SpecialStringTok{/v1/audio/translate{-}bark"}\NormalTok{,}
\NormalTok{            headers}\OperatorTok{=}\NormalTok{\{}\StringTok{"Authorization"}\NormalTok{: }\SpecialStringTok{f"Bearer }\SpecialCharTok{\{}\NormalTok{token}\SpecialCharTok{\}}\SpecialStringTok{"}\NormalTok{\},}
\NormalTok{            data}\OperatorTok{=}\NormalTok{\{}\StringTok{"pet\_id"}\NormalTok{: args.pet\_id\},}
\NormalTok{            files}\OperatorTok{=}\NormalTok{\{}\StringTok{"audio"}\NormalTok{: (args.audio.name, recording, }\StringTok{"audio/wav"}\NormalTok{)\},}
\NormalTok{            timeout}\OperatorTok{=}\DecValTok{30}\NormalTok{,}
\NormalTok{        )}
    \BuiltInTok{print}\NormalTok{(response.status\_code, response.text)}
\NormalTok{    response.raise\_for\_status()}


\ControlFlowTok{if} \VariableTok{\_\_name\_\_} \OperatorTok{==} \StringTok{"\_\_main\_\_"}\NormalTok{:}
\NormalTok{    main()}
\end{Highlighting}
\end{Shaded}

\subsubsection{Webhook
бүртгүүлэх}\label{webhook-ux431ux4afux440ux442ux433ux4afux4afux43bux44dux445}

\texttt{subscribe\_webhook.py} нь pet ID ба өөрийн эзэмшлийн HTTPS
callback URL-ийг JSON-аар илгээнэ. Callback endpoint-ийг олон нийтэд
ажиллуулж туршаагүй; \texttt{url} field-ийн нэрийг мөн бодит
contract-той шалгах хэрэгтэй.

\begin{Shaded}
\begin{Highlighting}[]
\CommentTok{"""UE{-}5 Corg.ly exercise: subscribe to pet events. Contract is unverified."""}

\ImportTok{import}\NormalTok{ argparse}
\ImportTok{import}\NormalTok{ os}
\ImportTok{from}\NormalTok{ urllib.parse }\ImportTok{import}\NormalTok{ urlparse}

\ImportTok{import}\NormalTok{ requests}


\KeywordTok{def}\NormalTok{ main() }\OperatorTok{{-}\textgreater{}} \VariableTok{None}\NormalTok{:}
\NormalTok{    parser }\OperatorTok{=}\NormalTok{ argparse.ArgumentParser(description}\OperatorTok{=}\StringTok{"Register an HTTPS webhook with a lab server"}\NormalTok{)}
\NormalTok{    parser.add\_argument(}\StringTok{"webhook\_url"}\NormalTok{, }\BuiltInTok{help}\OperatorTok{=}\StringTok{"HTTPS callback owned by the caller"}\NormalTok{)}
\NormalTok{    parser.add\_argument(}\StringTok{"{-}{-}pet{-}id"}\NormalTok{, required}\OperatorTok{=}\VariableTok{True}\NormalTok{, }\BuiltInTok{help}\OperatorTok{=}\StringTok{"Registered pet ID"}\NormalTok{)}
\NormalTok{    args }\OperatorTok{=}\NormalTok{ parser.parse\_args()}
    \ControlFlowTok{if}\NormalTok{ urlparse(args.webhook\_url).scheme }\OperatorTok{!=} \StringTok{"https"}\NormalTok{:}
\NormalTok{        parser.error(}\StringTok{"webhook\_url must use HTTPS"}\NormalTok{)}
\NormalTok{    base\_url }\OperatorTok{=}\NormalTok{ os.environ[}\StringTok{"CORGLY\_BASE\_URL"}\NormalTok{].rstrip(}\StringTok{"/"}\NormalTok{)}
\NormalTok{    token }\OperatorTok{=}\NormalTok{ os.environ[}\StringTok{"CORGLY\_TOKEN"}\NormalTok{]}
\NormalTok{    response }\OperatorTok{=}\NormalTok{ requests.post(}
        \SpecialStringTok{f"}\SpecialCharTok{\{}\NormalTok{base\_url}\SpecialCharTok{\}}\SpecialStringTok{/v1/webhooks/subscribe"}\NormalTok{,}
\NormalTok{        headers}\OperatorTok{=}\NormalTok{\{}\StringTok{"Authorization"}\NormalTok{: }\SpecialStringTok{f"Bearer }\SpecialCharTok{\{}\NormalTok{token}\SpecialCharTok{\}}\SpecialStringTok{"}\NormalTok{\},}
\NormalTok{        json}\OperatorTok{=}\NormalTok{\{}\StringTok{"pet\_id"}\NormalTok{: args.pet\_id, }\StringTok{"url"}\NormalTok{: args.webhook\_url\},}
\NormalTok{        timeout}\OperatorTok{=}\DecValTok{20}\NormalTok{,}
\NormalTok{    )}
    \BuiltInTok{print}\NormalTok{(response.status\_code, response.text)}
\NormalTok{    response.raise\_for\_status()}


\ControlFlowTok{if} \VariableTok{\_\_name\_\_} \OperatorTok{==} \StringTok{"\_\_main\_\_"}\NormalTok{:}
\NormalTok{    main()}
\end{Highlighting}
\end{Shaded}

\subsection{Code sample audit}\label{code-sample-audit}

Оноо 1--5; хамгийн бага нотолгоотой Trustworthy хэмжүүрийг өндөр
оноогоор нөхөөгүй. Bhatti х.89--94-ийн request/response-ийн яг тохирол,
товч, ойлгомжтой, солих утга тодорхой байх шалгуурыг хэрэглэв.

\begingroup\footnotesize\setlength{\tabcolsep}{2pt}

\begin{longtable}[]{@{}lrrrrrr@{}}
\toprule\noalign{}
Sample & Explained & Concise & Clear & Usable & Trustworthy & Дундаж \\
\midrule\noalign{}
\endhead
\bottomrule\noalign{}
\endlastfoot
Photo upload & 5 & 4 & 5 & 3 & 1 & 3.6 \\
Bark translation & 5 & 4 & 5 & 3 & 1 & 3.6 \\
Webhook subscribe & 5 & 4 & 5 & 3 & 1 & 3.6 \\
Nogoolin OpenAPI request/response examples & 4 & 5 & 4 & 3 & 2 & 3.6 \\
\textbf{Нийт} & & & & & & \textbf{3.6/5} \\
\end{longtable}

\endgroup

Өндөр үнэлгээний үндэслэл: зорилго, оролт, тайлбартай нэр, богино API
call бий; \texttt{foo/bar/test} орлуулагч байхгүй. Бага үнэлгээний
үндэслэл: Corg.ly request field болон response shape бодитоор
батлагдаагүй, Nogoolin example JSON нь код/схемээр нягталсан загвар
боловч live capture биш. Энэ нь DoD-ийн 4/5 босгыг \textbf{одоогоор
хангахгүй}. Бодит sandbox ирвэл field contract-ийг тулгаж, response-ийг
capture хийж, дахин үнэлнэ.

\subsection{Swagger UI ба Redoc --- шийдвэрийн тайлан (100
үг)}\label{swagger-ui-ux431ux430-redoc-ux448ux438ux439ux434ux432ux44dux440ux438ux439ux43d-ux442ux430ux439ux43bux430ux43d-100-ux4afux433}

Swagger UI нь endpoint бүрийн параметр, хамгаалалт, body-г бөглөж хүсэлт
илгээх боломжтой тул хөгжүүлэгчийн локал туршилтад тохирно. Гэхдээ ``Try
it out'' нь токен, ажиллаж буй API, өгөгдлийн сан шаардана; нийтэд
нээлттэй орчинд санамсаргүй бичилт үүсгэж болно. Redoc нь уншихад
цэгцтэй navigation, schema холбоос, гурван баганат үзэмжтэй тул
хэрэглэгчид API гэрээг тайлбарлахад илүү тохиромжтой. Хайлт ба байршлын
хурдыг endpoint дээр хэмжээгүй. Одоохондоо Redoc-ийг унших үндсэн
хувилбар, Swagger UI-г зөвхөн локал sandbox гэж сонгов. Хоёр
renderer-ийг ижил YAML-аас үүсгэвэл зөрүү багасна. Public URL,
хэрэглэгчийн үнэлгээ байхгүй тул энэ шийдвэрийг түр гэж үзэж,
deployment-ийн дараа дахин шалгана. Энэ нь лабораторийн хүрээний санал
юм.

\subsection{Reflection}\label{reflection}

Bhatti Ch.5 х.89--94-ийн хамгийн чухал санаа нь response жишээ код шигээ
шалгагдах ёстой гэдэг байлаа. Өмнөх spec-д жишээ өгөгдөл байхад түүнийг
бодит үр дүн гэж андуурах амархан. Би энэ удаа жишээг зориуд
``illustrative'' гэж тэмдэглэж, өөрийн хийж чадаагүй live шалгалтыг ил
үлдээв. Chinchilla Ch.6 х.79-ийн docs-as-code ба render-ийн санаагаар
YAML-г локал lint, standalone HTML, PDF тайлантай нэг эхээс холбоно.
Миний төслийн хамгийн том зөрүү нь хуучин Phase-0 spec-ийн зарим
response жишээ бодит controller-ийн return envelope-той бүрэн нийцэхгүй
байгаа явдал; W5-ийн таван route-ыг кодоос дахин тулгасан ч бүх API-ийн
drift audit хийгээгүй.

\subsection{Нотолгоо ба дутуу
шалгуур}\label{ux43dux43eux442ux43eux43bux433ux43eux43e-ux431ux430-ux434ux443ux442ux443ux443-ux448ux430ux43bux433ux443ux443ux440}

\begin{itemize}
\tightlist
\item
  \texttt{redocly\ lint}: OpenAPI valid, алдаа 0; localhost болон лиценз
  metadata-ийн warning-ийг README-д тэмдэглэв.
\item
  \texttt{python3\ -m\ py\_compile}: гурван script синтаксийн хувьд зөв.
\item
  Redoc standalone HTML болон Swagger UI local HTML-г бэлдэв. Static
  asset замыг шалгасан ч browser rendering, ``Try it out'' runtime
  туршилт хийгдээгүй.
\item
  Public хоёр URL, workspace upload, Git commit/push, CI executable
  bonus, Corg.ly live response: хийгдээгүй.
\end{itemize}

Эх:
\href{Sprint\%2005\%20Lab\%20Instructions\%20-\%20OpenAPI\%203.0\%20\%26\%20Bhatti\%20Code\%20Samples.pdf}{W5
тусгай handout}; Bhatti et al., \emph{Docs for Developers}, Ch.5
х.86--94; Chinchilla, \emph{Technical Writing for Software Developers},
Ch.6 х.79.

\end{document}
